El presente proyecto propone la implementación de un sistema de 
interacción adaptativa en realidad virtual (VR, por sus siglas 
en inglés) para una experiencia de videoclip gamificado. Su 
enfoque es el desarrollo de un mecanismo capaz de registrar 
métricas de desempeño del usuario en tiempo real, interpretar 
estados de interacción y ajustar dinámicamente parámetros del 
\textit{gameplay}, como la velocidad, cantidad de objetos y complejidad 
de las tareas. Esto, basándose en el diseño de interfaces de 
interacción en VR, técnicas de interacción en entornos 
inmersivos, adaptación basada en experiencia del jugador, 
retroalimentación multisensorial y evaluación de presencia 
y experiencia de usuario.

Para el desarrollo del sistema se utiliza Unity para la lógica 
interactiva y el prototipado de la experiencia, mientras que los 
modelos en 3D se crean en Blender. El desarrollo conlleva un 
prototipo funcional que integre entrada de VR, detección de 
interacción, almacenamiento de métricas en tiempo real, lógica de 
estados del jugador y adaptación dinámica. Como resultado, se 
busca aportar una base técnica y metodológica para experiencias 
inmersivas que no solo respondan a las entradas del usuario, sino
también ajusten su comportamiento en función del desempeño y de 
indicadores de involucramiento durante la experiencia.
